\SourceLeaf{193}\clearpage
\begin{center}
{\Large\textsc{Nachlass und Briefwechsel.}}
\end{center}

\bigskip
\begin{center}[ZUR PARALLELENTHEORIE.]\end{center}
\begin{center}\rule{2.5cm}{0.4pt}\end{center}

\begin{center}
\textsc{Gauss} an \textsc{Taurinus}. Göttingen, 8. November 1824.
\end{center}

\begin{center}Ewr. Wohlgeboren\end{center}

gefälliges Schreiben vom 30. Oct. nebst dem beigefügten kleinen Aufsatz habe ich nicht ohne
Vergnügen gelesen, um so mehr, da ich sonst gewohnt bin, bei der Mehrzahl der Personen, die
neue Versuche über die sogenannte Theorie der Parallellinien [machen,] gar keine Spur von
wahrem geometrischen Geiste anzutreffen. Gegen Ihren Versuch habe ich nichts (oder nicht viel)
anderes zu erinnern, als dass er unvollständig ist. Zwar lässt Ihre Darstellung des Beweises,
dass die Summe der drei Winkel eines ebenen Dreiecks nicht grösser als $180^{0}$ sein kann, in
Rücksicht auf geometrische Schärfe noch zu desideriren übrig. Allein diess würde sich
ergänzen lassen, und es leidet keinen Zweifel, dass jene Unmöglichkeit sich auf das
allerstrengste beweisen lässt. Ganz anders verhält es sich aber mit dem $2^{\mathrm n}.$ Theil,
dass die Summe der Winkel nicht kleiner als $180^{0}$ sein kann; diess ist der eigentliche
Knoten, die Klippe, woran alles scheitert. Ich vermuthe, dass Sie Sich noch nicht lange
